\documentclass[10pt,landscape]{article}
\usepackage{multicol}
\usepackage[landscape, inner=1cm, outer=1cm, top=1cm, bottom=1cm]{geometry}
\usepackage{amsmath,amsthm,amsfonts,amssymb}
\usepackage{color,graphicx,overpic}
\usepackage{hyperref}
\usepackage{amsthm}
\usepackage{mathrsfs}
\usepackage{enumerate}
\usepackage{enumitem}
\usepackage{icomma}
\usepackage{soul}
\usepackage{physics}
\usepackage{cancel}
\usepackage[european]{circuitikz}

% Turn off header and footer
\pagestyle{empty}

% Redefine section commands to use less space
\makeatletter
\renewcommand{\section}{\@startsection{section}{1}{0mm}%
	{-1ex plus -.5ex minus -.2ex}%
	{0.5ex plus .2ex}%x
	{\normalfont\large\bfseries}}
\renewcommand{\subsection}{\@startsection{subsection}{2}{0mm}%
	{-1explus -.5ex minus -.2ex}%
	{0.5ex plus .2ex}%
	{\normalfont\normalsize\bfseries}}
\renewcommand{\subsubsection}{\@startsection{subsubsection}{3}{0mm}%
	{-1ex plus -.5ex minus -.2ex}%
	{1ex plus .2ex}%
	{\normalfont\small\bfseries}}
\makeatother

% Don't print section numbers
\setcounter{secnumdepth}{0}

\setlength{\parindent}{0pt}
\setlength{\parskip}{0pt plus 0.5ex}

% My environments
\theoremstyle{definition}
\newtheorem*{primer}{Primer}

\begin{document}
\begin{multicols}{3}

  \section{Senzorji}

  Imamo vhodne podatke \( z(t) \) v senzor in izhodne podatke izven \( x(t) = \hat{x} (t) \). Senzor \( n \)-tega reda je optimalen za sledenje spremenljivki z dinamiko \( \frac{\mathrm{d} ^{n}}{\mathrm{d} t^{n}} x = w(t) \), kjer je \( w(t) \) dinamičen šum. Ojačevalni faktor je \( \lim_{t \to \infty} K = K_{\infty} \). 

  \subsection{Senzor 1. reda}

  Enačba senzorja je \( z(t) = \tau \dot{x} + x \), kjer je \( \tau ^2 = \frac{R}{Q} \). \( R = \left\langle r(t), r(t') \right\rangle \) in \( Q = \left\langle w(t) w(t') \right\rangle \). Prenosna funkcija \( H(s) = \frac{1}{1 + \tau s} \) in potem izračunamo \( X(s) = H(s) Z(s) \). 

  \begin{primer}[Linearen signal] Signal \( z(t) = kt \), kar je \( Z(s) = \frac{k}{s ^2} \). Izhodni signal je (parcialni ulomki) \\
    \( X(s) = \frac{1}{1 + \tau s} \frac{k}{ s^2 } = k \left( \frac{\tau ^2}{1 + \tau s} - \frac{\tau}{s} + \frac{1}{s ^2} \right) \). \\
    Izhodni signal je potem \\
    \( x(t) = k \left( \tau e^{-t/ \tau} - \tau \Theta(t) + t \right) \) in v \( \lim_{t \to \infty} = k \tau - kt \). Potem je \\
    \( \Delta (t) = z(t) - k(t) = k \tau \)
  \end{primer}

  \subsection{Senzor 2. reda}
  
  Enačba senzorja je \( \ddot{x} + 2 \xi \omega_0 \dot{x} + \omega_0 ^2 x = \omega_0 ^2 z \), kjer je \( \omega_0 ^2 = \frac{Q}{R} \) in optimalna frekvenca dušenja \( \theta = \frac{\sqrt{2}}{2} \). Prenosna funkcija je \( H(s) = \left( \omega_0 ^2 + \color{gray}{2 \xi \omega_0 s} \right) \left[ s ^2 + 2 \xi \omega_0 s + \omega_0 ^2 \right]^{-1} = \omega_0 ^2 \left[ \left( s + \xi \omega_0 \right)^2 + \omega_0 ^2 \left( 1 - \xi \right) \right] ^{-1} \). Sivi del je za sledenje še odvodu \( \dot{z} \).

  Tri vrste dušenja: \( \xi < 1 \) je podkritično, \( \xi > 1 \) nadkritično in \( \xi = 1 \) kritično. 

  \subsection{Periodični vhodni signal}

  \( z(t) = z_0 e^{\mathrm{i} \omega t} \implies x(t) = x_0 e^{\mathrm{i} \delta} e^{\mathrm{i} \omega t} \)  in \( s \to \mathrm{i} \omega \).

  Velikosti amplitud \( \left| H(\mathrm{i} \omega) \right| = \frac{x_0}{z_0}\) in \( \tan \delta = \frac{\Im H}{\Re H} \).
  
  \section{Elektronika}

  Tok ne teče v smeri \( U_{out} \), ampak na ground. 

  \textbf{low-pass filter}

  $H = \frac{U_{out}}{U_{in}} = \frac{1}{1 + i \omega \tau}$, \qquad $\tau = RC$ \\
  \vspace{-15pt}
  \begin{center}
    \begin{circuitikz}
      \scalebox{0.6}{
        \draw node[left]{$U_{in}$} (0, 0) to[R=$R$, o-] (2, 0) to[short, -o] (4, 0) 
        node[right]{$U_{out}$};
        \draw (2, 0) to[C=$C$] node[ground]{} (2, -1.5);
      }
    \end{circuitikz} 
  \end{center}\vspace{-25pt}

  Bodejev diagram ima na \( x \) osi \( \log_{10} (\omega t) \) in na \( y \) osi \( 10 \log \left( \left| H \right| \right) \).

  \[
    \begin{array}[tb]{|c|c|c|} \hline
      \omega \tau \gg 1 & \left( \omega \tau \right)^{-2} & -20 \log (\omega \tau) \\
      \omega \tau \ll 1 & 1 & 0 \mathrm{dB} \\
      \omega \tau = 1 & \frac{1}{2} & 10 \log 2^{-1} = - 3 \, \mathrm{dB}. \\\hline
    \end{array}
  \]

  \textbf{high-pass filter}

  $H = \frac{i \omega \tau}{1 + i \omega \tau}$, \qquad $\tau = RC$ \\
  \vspace{-20pt}
  \begin{center}
    \begin{circuitikz}
      \scalebox{0.6}{
        \draw node[left]{$U_{in}$} (0, 0) to[C=$C$, o-] (2, 0) to[short, -o] (4, 0) 
        node[right]{$U_{out}$};
        \draw (2, 0) to[R=$R$] node[ground]{} (2, -1.5);
      }
    \end{circuitikz}
  \end{center}\vspace{-30pt}

  \section{Aktivna vezja}
  \textbf{Operacijski ojačevalec (OPAMP)} \\
  $U_{out} = A(s) (U_+ - U_-)$ \\
  Idealni OPAMP: $A_s \to \infty$, $I_+ = I_- = 0$ in \( Z_+ = Z_- = \infty \) \\
  \vspace{-60pt}
  \begin{center}
    \begin{circuitikz}
      \scalebox{0.6}{
        \draw (0,0) node[op amp] (opamp) {} 
        (opamp.+) node[left] {$U_+$}
        (opamp.-) node[left] {$U_-$}
        (opamp.out) node[right] {$U_{out}$}
        (opamp.up) --++(0,0.5) node[vcc]{$5 \ {V}$}
        (opamp.down) --++(0,-0.5) node[vee]{$-5 \ {V}$};
      }
    \end{circuitikz} 
  \end{center}\vspace{-40pt}


  \textbf{Diferenciator} \\
  $U_{out}(t) = \tau \dot U_{in}(t)$ \\
  $H = i \omega \tau$, \qquad $\tau = RC$ \\

  V stičišču pred opampom in po kondenzatorjem, je napetost \( 0 \mathrm{V} \), saj je \( U_+ = U_- \) in \( U_- = 0 \). Tok teče od \( U_{in} \) do \( U_{out} \) in skozi člena teče enak tok \( I_c = I_R \), kjer je \( I_c = U_{in} / Z_C \) in \( I_R = U_{out} / R \)
  \vspace{-25pt}
  \begin{center}
  \begin{circuitikz}
    \scalebox{0.6}{
      \draw (0,0) node[op amp] (opamp) {} 
      (opamp.+) to[short] ++(0, -0.5) node[ground] {}
      (opamp.-) to[short] ++(-0.5, 0) --coordinate(A) ++(-0.5, 0) to[C=$C$] ++(-1, 0) to[short, 
      -o] ++(-1, 0) node[left]{$U_{in}$}
      (A) to[short] ++(0, 1) to[R=$R$] ++(3, 0) -- ++(0, -1.49)
      (opamp.out) to[short, -o] (2, 0) node[right] {$U_{out}$};
    }
  \end{circuitikz}
  \end{center}\vspace{-25pt}

  \textbf{Integrator} \\
  $U_{out}(t) = - \frac{1}{\tau} \int_0^t U_{in}(t') \dd{t'}$ \\
  $H = - \frac{1}{i \omega \tau}$, \qquad $\tau = RC$ \\
  \vspace{-25pt}
  \begin{center}
  \begin{circuitikz}
    \scalebox{0.6}{
      \draw (0,0) node[op amp] (opamp) {} 
      (opamp.+) to[short] ++(0, -0.5) node[ground] {}
      (opamp.-) to[short] ++(-0.5, 0) --coordinate(A) ++(-0.5, 0) to[R=$R$] ++(-1, 0) to[short, 
      -o] ++(-1, 0) node[left]{$U_{in}$}
      (A) to[short] ++(0, 1) to[C=$C$] ++(3, 0) -- ++(0, -1.49)
      (opamp.out) to[short, -o] (2, 0) node[right] {$U_{out}$};
    }
  \end{circuitikz}
  \end{center}\vspace{-15pt}

  \textbf{Invertirajoči ojačevalnik} \\
  $H = - \frac{R_2}{R_1}$ \\
  \vspace{-30pt}
  \begin{center}
  \begin{circuitikz}
    \scalebox{0.6}{
      \draw (0,0) node[op amp] (opamp) {} 
      (opamp.+) to[short] ++(0, -0.5) node[ground] {}
      (opamp.-) to[short] ++(-0.5, 0) --coordinate(A) ++(-0.5, 0) to[R=$R_1$] ++(-1, 0) to[short, 
      -o] ++(-1, 0) node[left]{$U_{in}$}
      (A) to[short] ++(0, 1) to[R=$R_2$] ++(3, 0) -- ++(0, -1.49)
      (opamp.out) to[short, -o] (2, 0) node[right] {$U_{out}$};
    }
  \end{circuitikz}
  \end{center}\vspace{-25pt}

  \textbf{Neinvertirajoči ojačevalnik} \\

  $H = 1 + \frac{R_1}{R_2}$ \\
  \vspace{-60pt}
  \begin{center}
  \begin{circuitikz}
    \scalebox{0.6}{
      \draw (0,0) node[op amp] (opamp) {} 
      (opamp.+) to[short, -o] ++(-1,0) node[left]{$U_{in}$}
      (opamp.-) -- ++(0, 1) -- ++(2.5, 0) coordinate(A) to[R=$R_1$] ++(0, -1.49)
      (A) to[R=$R_2$] ++(2, 0) node[ground]{}
      (opamp.out) to[short, -o] (2, 0) node[right] {$U_{out}$};
    }
  \end{circuitikz}
  \end{center}

  \textbf{Napetostni sledilnik} \\ 

  Po dolgem času je \( U_{in} = U_{out} \) in prenosna funkcija je \( H = 1 \). Pomaga nam pri združevanje vezij, saj v negativno zanko ne teče tok!
  \vspace{-35pt}
  \begin{circuitikz}
    \scalebox{0.6}{
      \draw (0,0) node[op amp] (opamp) {} 
      (opamp.+) to[short, -o] ++(-1,0) node[left]{$U_{in}$}
      (opamp.-) -- ++(0, 1) -- ++(2.5, 0) coordinate(A) -- ++(0, -1.49)
      (opamp.out) to[short, -o] (2, 0) node[right] {$U_{out}$};
    }
  \end{circuitikz} \vspace{20pt}\\

  \textbf{Negativna povratna zanka}

  Za prenosno funkcijo \( H(s) \) in ojačevalnik \( A \), in dodatno prenosno funkcijo \( F(s) \) bo veljalo \( U_{out} = K H \left( U_{in} - F(s) U_{out} \right) \).

  \subsection{Resonančna vezja}

  \textbf{Pasovni prepustni (band-pass) filter} \\
  $H = \frac{s \omega_c}{s^2 + s \omega_c + \omega_0^2}$\\
  $\omega_0^2 = \frac{1}{LC}$, \qquad $\omega_c = \frac{1}{RC}$ \\
  Blizu resonance $s \approx i \omega_0$: \\
  $\Delta \omega_{-3 \, \text{dB}} = \omega_c$ \\
  $Q = \frac{\omega_0}{\Delta \omega_{-3 \, \text{dB}}} = \frac{\omega_0}{\omega_c}$ \\

  Rešujemo ga tako, da vzporedno vezana \( L \) in \( C \) nadomestimo z impedanco \( Z_{LC} \), za katero velja \( Z_{LC}^{-1} = Z_C^{-1} + Z_L^{-1} \), kjer je \( Z_C = \left( \mathrm{i} \omega C \right)^{-1} \) in \( Z_L = \mathrm{i} \omega L \).
  \vspace{-17pt}
  \begin{center}
    \begin{circuitikz}
      \scalebox{0.6}{
        \draw node[left]{$U_{in}$} (0, 0) to[R=$R$, o-] (2, 0) to[short, -o] (4, 0) 
        node[right]{$U_{out}$};
        \draw (2, 0) -- ++(0, -0.5) -- ++(0.5, 0) to[cute, L=$L$] ++(0, -1.5) -- 
        ++(-0.5, 
        0) 
        coordinate(A) -- ++(-0.5, 0) to[C=$C$] ++(0, 1.5) -- ++(0.5, 0)
        (A) node[ground]{};
      }
    \end{circuitikz} 
  \end{center}\vspace{-35pt}

  \textbf{Pasovni neprepustni (band-reject) filter} \\
  $H = \frac{s^2 + \omega_0^2}{s^2 + s\omega_c + \omega_0^2}$ \\
  $\omega_0^2 = \frac{1}{LC}$, \qquad $\omega_c = \frac{R}{L}$ \\
  Blizu resonance $s \approx i \omega_0$: \\
  $\Delta \omega_{-3 \, \text{dB}} = \omega_c$ \\
  $Q = \frac{\omega_0}{\Delta \omega_{-3 \, \text{dB}}} = \frac{\omega_0}{\omega_c}$ \\
  \vspace{-17pt}
  \begin{center}
    \begin{circuitikz}
      \scalebox{0.6}{
        \draw node[left]{$U_{in}$} (0, 0) to[R=$R$, o-] (2, 0) to[short, -o] (4, 0) 
        node[right]{$U_{out}$};
        \draw (2, 0) to[C=$C$] ++(0, -1) to[cute, L = $L$] ++(0, -1) node[ground]{};
      }
    \end{circuitikz} 
  \end{center}\vspace{-45pt}

  \section{Statistika}

  Vzorec \( \left\{ z_i \right\}_{i = 1}^N \) in \( z_i \sim \mathcal{N} (a, \sigma ^2 ) \). Želimo oceniti vrednosti parametrov \( \mathcal{N} \). Pogoj nepristranskosti pravi, da je pričakovana vrednosti cenilke \( u \) enaka količini, ki jo ocenjujemo: \( \left\langle u \right\rangle = a\). \\

  \( \bar{z} = \frac{1}{N} \sum\limits_i^{} z_i \) \\

  Če poznamo \( a \),  \( \sigma = \frac{1}{N} \sum\limits_{}^{} \left( z_i - a \right) ^2 \), drugače \\

  \( \left\langle s ^2 \right\rangle = \frac{1}{N - 1} \sum\limits_i^{} \left( z_i - \bar{z} ^2 \right)\)\\
  \textbf{T-statistika} \\
  Porazdelitev \( S(k) = \frac{\mathrm{d} P}{\mathrm{d} T} (k) = \frac{1}{\sqrt{k} B \left( \frac{k}{2}, \frac{1}{2} \right)} \left( 1 + \frac{T ^2}{k} \right)^{- \frac{k + 1}{2}} \) in \\
  \( T = \frac{\bar{z} - a}{\sqrt{s ^2}} \sqrt{N} \sim S (N - 1) \), \( k = N - 1 \) prostostne stopnje \\
  Interval zaupanja pir \( \alpha \) določiš z rešitvijo \( P \left( \left| T \right| > T_{>} \right) = \bf{\alpha} \) in ocena je \\
  \( a = \bar{z} \pm T_{>} \frac{s}{\sqrt{N}} \)

  \( \chi ^2 \) \textbf{statistika} \\
  Izmerke \( z_i \) pretvorimo na nromalno porazdelitev \( x_i = \frac{\left( z_i - a \right) ^2}{\sigma ^2} \sim \mathcal{N} (0, 1) \), potem \\
  \( \chi ^2 = \sum\limits_i^{} \frac{\left( z_i - a \right) ^2}{\sigma ^2} = (N - 1) \frac{s ^2}{\sigma ^2} \sim \chi ^2 (k) = \chi ^2 (N - 1 \), \\
  kjer je \( \chi ^2 \) porazdelitev \( \chi ^2 (k) = \frac{\mathrm{d} P}{\mathrm{d} \chi ^2 }(k) \). \\
  Interval zaupanja pri \( \alpha \), rešimo enačbi \( P \left( \chi ^2 > \chi_{>} ^2 \right) = \frac{\alpha}{2} \) in \( P \left( \chi ^2 > \chi_{<} ^2 \right) = 1 - \frac{\alpha}{2} \)

  \( \sigma_{\lessgtr} = (N - 1) \frac{s ^2}{\chi ^2_{\gtrless}} \)

  \textbf{Primerjava dveh vzorcev} \\
  Dva vzorca \( \left\{ x_i \right\}_{i = 1}^N \) in \( \left\{ y_i \right\}_{i = 1}^M \), \( x_i \sim \mathcal{N} (a_x, \sigma ^2 \) in \( y\sim \mathcal{N} ( a_y, \sigma ^2 ) \). 

  \[ T = \frac{\left( \bar{x} - \bar{y} \right) - \left( a_x - a_y \right)}{\sqrt{\frac{1}{N} + \frac{1}{M}}} \frac{\sqrt{N + M - 2}}{\sqrt{s_x ^2 (N - 1) + s_y ^2 (M - 1}} \sim S (N + M - 2) 
  \]

  \textbf{Testiranje hipotez}\\
  Nulta hipoteza \( H_0: \ a = a_0 \). Izračunaš \( T_0 = \frac{\bar{z} - a}{s ^2} \sqrt{N} \) in primerjaš z rešitvijo enačbe \( P \left( \left| T \right| > T_C \right) = \alpha \). Če \( \left| T_0 \right| > T_C \) hipotezo zavrnemo, drugače je ne moremo zavreči.

  Nulta hipoteza \( H_0 = \sigma ^2 = \sigma ^2_0 \). Izračunaš \( \chi_0 ^2 = (N - 1) \frac{s ^2}{\sigma_0 ^2} \) in primerjaš z rešitvama \\
  \( P \left( \chi ^2 > \chi_{>} ^2 \right) = \frac{\alpha}{2} \) in \( P \left( \chi ^2 > \chi_{<} ^2 \right) = 1 - \frac{\alpha}{2} \)\\
  Če \( \chi_0 ^2 \le \chi_{<} ^2 \) ali \( \chi_0 ^2 \ge \chi_{>} ^2 \), hipotezo zavrnemo. 

  \section{Oblikovni testi}

  \subsection{Pearsonov test \( \chi_p ^2 \)}

  Izmeri \( \left\{ z_i \right\}_{i = 1}^N \) so porazdeljeni v \( k \) razredov/intervalo. Število izmerkov v \( k \)-tem intervalu je \( N_k \) in testiramo porazdelitveno funkcijo \( \frac{\mathrm{d} P}{\mathrm{d} x} \) \\

  \( P_k = \int\limits_{z_{k - 1}}^{z_k} \frac{\mathrm{d} P}{\mathrm{d} x} \, \mathrm{d} x \) \\
  \( \chi_p ^2 = \sum\limits_{i = 1}^k \frac{(N_k - Np_k)}{N p_k} \sim \chi ^2 (k - 1) \) \textbf{za} \( \bf{Np_k} \ge 5 \) . 

  \section{Linearna metoda najmanjših kvadratov}

  Imamo \( N \) meritev \( z_i = z(t_i) \sim \mathcal{N} (a_i, \sigma_i ^2), \ \sigma_{ij} = 0 \) in iščemo optimalno oceno \( M \) parametrov \( \mathbf{x} \in \mathbb{R}^M \) za model

  \[ \mathbf{z} = \begin{bmatrix} z(t_{1}) \\ z(t_{2}) \\ \vdots \\ z(t_{N}) \end{bmatrix} = H \mathbf{x} = \begin{bmatrix}
    f_{0} (t_1) & \ldots & f_M (t_1)  \\
    \vdots & & \\
    f_0 (t_{N}) & \ldots & f_M (t_N)
  \end{bmatrix}
  \begin{bmatrix} x_{0} \\ x_{2} \\ \vdots \\ x_{M} \end{bmatrix}
  \]
Ker je \( \sigma = \sigma_i \ \forall i  \), je \( R = \sigma ^2 I \) (ni korelacije).

  Minimiziranje kvadratne forme \\
  \( \chi ^2 = 2J = \mathbf{r} R^{-1} \mathbf{r}^T \sim \chi ^2 (N - 1) \),\\
  poda enačbi\\
  \( \hat{\mathbf{x}} = \left( H^T H \right)^{-1} H^T \mathbf{z} \) in \\
  \( P = \left( H^T H \right)^{-1} \delta ^2 \)

  \section{Linearni model z 2 parametroma}

  \( z = x_0 t + x_1 \), kar pomeni \( f_0 (x) = x \) in \( f_1 (x) = 1 \). Imamo \( n \) izmerkov \( z_i = z (t_i) \). Potem je
  \[ \hat{\mathbf{x}} = \frac{1}{C} \begin{bmatrix} N \sum\limits_i^{} t_i z_i  - \left( \sum\limits_i^{} t_i \right) \cdot \left( \sum\limits_{i}^{}z_i \right) \\ - \left( \sum\limits_i^{} t_i \right) \left( \sum\limits_i^{} t_i z_i \right) + \left( \sum\limits_i^{} z_i \right) \cdot \left( \sum\limits_i^{} t_i ^2 \right)  \end{bmatrix},
  \]
  kjer je \( C = N \sum\limits_i^{} t_i ^2 - \left( \sum\limits_i^{} t_i \right) ^2 \).

  \section{Misc}

  Laplaceove transformacije

  \[
    \begin{array}[tb]{|c|c|} \hline
      f(t) & F(s)  \\\hline
      \Theta(t) & \frac{1}{s}  \\
      e^{at} & \frac{1}{s - a} \\
      f(t) e^{at}& F(s - a) \\
      \frac{\mathrm{d} }{\mathrm{d} t} f(t) & s F(s) \\
      t & \frac{1}{s ^2} \\
      t^n & \frac{n!}{s^{n + 1}} \\
      \cos (\omega t) & \frac{s}{s ^2 + \omega ^2} \\
      \sin (\omega t) & \frac{\omega}{s ^2 + \omega ^2} \\
      \delta(t) & 1 \\
      \Theta(t - a) & \frac{1}{s} e^{-a s} \\
      f(t - \tau) \Theta(t - \tau) & e^{- s \tau} F(s) \\
      \frac{\mathrm{d} ^{n}}{\mathrm{d} t^n} f(t) & s^n F(s) \\ \hline
    \end{array}
  \]
  Pod pogojem, da so \( f(0) = f'(0) = \ldots = \frac{\mathrm{d} ^{n - 1}}{\mathrm{d} t^{n - 1}} f(0) = 0 \). Če \( f(0) \ne 0 \), potem \( \mathcal{L}(f) = F(s) + f(0) \).\\
  Inverz matrike \\
  \( A = \begin{bmatrix}
         a & b \\
         c & d
  \end{bmatrix} \) je\\
  \( A^{-1} = \frac{1}{ad - bc} \begin{bmatrix}
               d & -b \\
               -c & a
               \end{bmatrix} \)
\end{multicols}
\end{document}
